Our experimental design involves four main components: a sparse attention model, pathway metric extraction, EEG projection, and predictive validation.

\subsection*{Model and Data}

We use GPT-2 (124M parameters) as the base transformer. The model processes a fixed set of 10 prompts under 5 sparsity levels ($\{0.1, 0.3, 0.5, 0.7, 0.9\}$), generating 5 samples per condition (500 total samples, with 80/20 train-val split).

\subsection*{Experimental Protocol}

\begin{algorithm}
\caption{Causal EEG Hypothesis Testing Protocol}
\begin{algorithmic}[1]
\State \textbf{Input}: Fixed prompts, Sparsity levels $S$, Transformer $M$
\For{each sparsity level $s \in S$}
    \For{each prompt $p$}
        \State Set $M$.attention$\_$sparsity $\leftarrow s$
        \State Compute attention maps $A \leftarrow M(p)$
        \State Extract pathway metrics $\mathcal{P} \leftarrow \text{PathwayMetrics}(A)$
        \State Project to EEG: $\mathcal{E} \leftarrow \text{Projector}(\mathcal{P})$
        \State Record $(p, s, \mathcal{P}, \mathcal{E})$
    \EndFor
\EndFor
\State Train pathway predictor: $\hat{s}_{\mathcal{P}} = \text{MLP}_{\mathcal{P}}(\mathcal{P})$
\State Train EEG predictor: $\hat{s}_{\mathcal{E}} = \text{MLP}_{\mathcal{E}}(\mathcal{E})$
\State Evaluate transfer: $r_{transfer} = \text{corr}(\hat{s}_{\mathcal{E}}, s)$
\end{algorithmic}
\end{algorithm}

The key constraint is that sparsity is \textit{not directly encoded} in features; it acts only through attention distribution.
